\section{What we take from this}
\label{sec:discussion}

Each of the three questions ended up with a cleaner answer than we expected when we started.

For the bandit, the Lai--Robbins constant is a statement about slopes, not about ratios. On a three-armed Poisson problem even a million pulls is not asymptotic: KL-UCB approaches $C$ from above, Thompson sampling from below, and both have a large negative offset that makes $R_n/\log n$ look better than the bound. The other side of the bound is just as visible. UCB with $\sigma = 1$ has a beautiful median regret on every instance we tried, and it gets stuck in $38\%$ of the runs when the means are around 64. KL-UCB needs no tuning at all, and Thompson sampling was the best of the algorithms that never got stuck.

For the step sizes, the slowness of $1/n$ is a bias problem, not a noise problem. The exponent $\min(c(1-\gamma), 1/2)$ from stochastic approximation is matched to two decimals, and it says that $1/n$ forgets the initial values like $n^{-0.2}$ when $\gamma = 0.8$. In practice we would now use $c/(n+c-1)$ with $c$ around the horizon $1/(1-\gamma)$, and if $\gamma$ is not known, a polynomial step with $\omega$ around $0.6$ (or even $0.5$ for a fixed budget), keeping in mind that $\omega = 0.8$, which looks like a safe middle value, is bad for long horizons.

For the seeds, seven of the fourteen sentences we checked (thirteen from the report, one from the French follow-up) hold on every seed or almost, five hold on 80 to 89 seeds out of 100, and two were mostly produced by our particular seed. The fragile sentences are not random: they are the ones about small effects, about a final policy of a learner that never stops switching, or about the size of a bias that our seed happened to make extreme. If we rewrote the report today we would change exactly three things: drop the sentence on $Q(2, R)$ or give its value over many seeds, say that $\alpha = 0.5$ ends on the optimal policy about $88\%$ of the time, and replace the step $1/n$ by $5/(n+4)$ in the comparison of steps, since $1/n$ is not a fair representative of decreasing steps.

\paragraph{Limitations.} Everything here is on three toy problems where the answer is known, which is exactly what made the checks possible but also limits how far the conclusions go. The step-size experiments are only on policy evaluation of one policy of the robot. We did not try Q-learning with the $c/n$ rule beyond the quick test of the French follow-up, and we did not try iterate averaging (Polyak--Ruppert), which is the other classical way to get the $n^{-1/2}$ rate without knowing $\gamma$. For the bandits we used one family of Poisson instances and one horizon for the scale experiment. And 100 seeds is a lot for a course project but still gives intervals of a few percent on the claim frequencies.

\paragraph{What we would do next.} The most interesting open question for us is the second one in a setting where $\gamma$ is not known, or where the horizon changes from state to state. The heat map of Figure~\ref{fig:omega} suggests that a rule which adapts $\omega$ (or $c$) to the observed mixing time could work, and it would be nice to see if this can be done with a simple estimate. On the bandit side, it would be interesting to measure on the scale family how many pulls it takes before the local slope is within $10\%$ of $C$, as a function of $\mu$.

\paragraph{Reproducibility.} All scripts, raw results and figures are in the \texttt{research} folder of the repository. Each experiment runs in a few minutes on a laptop, the longest one (KL-UCB up to $10^6$) in about ten.
